 
\chapter{ Mediation Analysis: Natural Direct and Indirect Effects}
 \label{chapter::mediation}




当 treatment $Z$ 与 outcome $Y$ 之间存在一个中间变量 $M$ 时，由 $U=\{M(1),M(0)\}$ 定义的 principal strata 内的 causal effects 可以刻画 latent groups $U$ 之间的 treatment effect heterogeneity。若 $M$ 的确位于从 $Z$ 到 $Y$ 的 causal pathway 上，则某些 principal strata 内的 causal effects，例如 $\tau(1,1)$ 和 $\tau(0,0)$，可以提供关于 $Z$ 对 $Y$ 的 direct effect 的信息。不过，这些 direct effects 只针对两个 latent groups。另两个 principal strata 内的 causal effects，即 $\tau(1,0)$ 和 $\tau(0,1)$，同时包含 direct effects 与 indirect effects。从根本上说，principal stratification 并不提供关于 $Z$ 通过 $M$ 影响 $Y$ 的 indirect effect 的信息，因为它甚至不假设 $M$ 可以被干预。



在上面的讨论中，我以一种非正式的方式使用了 ``direct effect'' 和 ``indirect effect'' 这两个概念。当 $M$ 位于从 $Z$ 到 $Y$ 的 pathway 上时，研究者往往想评估 $Z$ 对 $Y$ 的效应有多大程度是通过 $M$ 传递的，又有多大程度是通过其它 pathways 传递的。这称为 {\it mediation analysis}。这正是本章的主题。



\section{Motivating Examples}





在 mediation analysis 中，我们有 treatment $Z$、outcome $Y$、mediator $M$，以及一些 background covariates $X$。\cref{fg::mediationDAG} 展示了它们之间的关系。下面给出几个具体例子。



\begin{figure}[ht]
\centering
$$
\begin{xy}
\xymatrix{
&& X \ar[dll] \ar[d] \ar[drr] \\
Z \ar[rr] \ar@/_1.5pc/[rrrr] & & M\ar[rr] & & Y}
\end{xy}
$$
\caption{Mediation analysis 的 causal diagram}\label{fg::mediationDAG}
\end{figure}


\begin{example}
\citet{Vanderweele::2012genetic} 使用 mediation analysis，评估 chromosome 15q25.1 上的 genetic variants 对 lung cancer 的效应有多大程度由 smoking 所 mediated，又有多大程度通过其它 causal pathways 起作用\footnote{回顾 \cref{chapter::evalue} 中 Fisher 关于一个 common genetic factor 同时影响 smoking behavior 与 lung cancer 的假说。}。Exposure levels 对应于 C alleles 数量从 0 到 2 的变化，smoking intensity 用每日吸烟支数的平方根度量，outcome 是 lung cancer indicator。\citet{Vanderweele::2012genetic} 的研究纳入了许多 sociodemographic covariates。  %case-control
\end{example}



\begin{example}
\citet{rudolph2018causal} 研究了从 neighborhood poverty 到 adolescent substance use 的 causal mechanism，其中 school and peer environment 作为 mediator。他们使用 National Comorbidity Survey Replication Adolescent Supplement 的数据；这是 2001--2004 年间进行的一项美国青少年全国代表性调查。Treatment 是 neighborhood disadvantage 的二元 indicator，其定义为：基于 2000 年 U.S. Census 数据，居住在 neighborhood socioeconomic status 最低三分位的社区。
四个二元 mediators 度量 school and peer environments，六个二元 outcomes 度量 substance use。
Baseline covariates 包括青少年的 sex、age、race、immigration generation、family income 等。
\end{example}


\begin{example}\label{eg::mediation-jobs}
\ri{R} 的 \ri{mediation} package 包含一个名为 \ri{jobs} 的数据集，它来自 JOBS II，一项研究 job training intervention 对 unemployed workers 效果的 randomized field experiment。
我们在 \cref{sec::cace-application} 中已经使用过这个数据集。
该项目不仅旨在提高失业者的 reemployment，也旨在改善 job seekers 的 mental health。因此，我们关心这个 intervention 通过 job search efficacy 对 mental health 产生的 indirect effect，以及它通过其它 pathways 起作用的 direct effect。稍后我们会在 \cref{section::BK-jobs} 中重新讨论这个例子。
\end{example}



\section{Nested Potential Outcomes}


\subsection{Natural Direct and Indirect Effects}
下面省略 unit $i$ 的下标 $i$，并假设所有 random variables 都是来自某个 superpopulation 的 IID draws。为简单起见，我们聚焦于 binary treatment $Z$。

 
我们首先考虑对 $z$ 的 hypothetical intervention，并定义与该 intervention 对应的 potential mediators 和 outcomes：
$$
\{ M(z) , Y(z): z=0,1\}.
$$
接着，我们考虑同时对 $z$ 和 $m$ 的 hypothetical intervention，并定义与这两个 interventions 对应的 potential outcomes：
$$
\{ Y(z, m): z=0, 1; m \in \mathcal{M}  \},
$$
其中 $ \mathcal{M} $ 包含 $m$ 的所有可能取值。
\citet{Robins::1992} 和 \citet{Pearl::2001} 进一步考虑与 intervention on $z$ 以及 $m=M(z')$ 对应的 nested potential outcomes：
$$
\{ Y(z, M_{z'}) : z = 0,1; z'=0,1   \}
$$
这里我把 $M(z')$ 写作 $M_{z'}$，以避免过多括号。记号 $Y(z,M_{z'})$ 表示如下 hypothetical outcome：treatment 被设为 level $z$，而 mediator 被设为它在 treatment $z'$ 下的 potential level $M(z')$。重要的是，$z$ 与 $z'$ 可以不同。
对于 binary treatment，总共有四个 nested potential outcomes：
$$
\{ Y( 1, M_1),    Y( 1, M_0),  Y( 0, M_1),  Y( 0, M_0) \}.
$$
Nested potential outcome $Y(1,M_1)$ 是如下 hypothetical outcome：treatment 被设为 $z=1$，mediator 被设为在 $z=1$ 下本来会取到的值。类似地，$Y(0,M_0)$ 是 treatment 被设为 $z=0$，mediator 被设为在 $z=0$ 下本来会取到的值时的 outcome。若 $Y(1,M_1)\neq Y(1)$ 或 $Y(0,M_0)\neq Y(0)$，那会很反直觉。
因此，本章始终采用下面的 composition assumption。

\begin{assumption}
[composition]\label{assume::composition}
$Y(z, M_z) =   Y(z)$ for $z=0,1$. 
\end{assumption}


Composition assumption 无法被证明。它确实是一个 assumption。若不引发哲学争论，我们甚至可以把 $Y(1)$ 定义为 $Y(1,M_1)$，把 $Y(0)$ 定义为 $Y(0,M_0)$。这样，\cref{assume::composition} 按定义成立。



Nested potential outcome $Y(1,M_0)$ 是如下 hypothetical outcome：unit 接受 treatment $1$，但它的 mediator 被设为没有 treatment 时的 natural value $M_0$。类似地，$Y(0,M_1)$ 是 unit 接受 control $0$，但它的 mediator 被设为 treatment 下的 natural value $M_1$ 时的 hypothetical outcome。它们是两个 cross-world counterfactual terms，可用于定义 direct and indirect effects。


\begin{definition}
[total, direct and indirect effects]
\label{def::mediation}
定义 treatment 对 outcome 的 total effect 为
$$
\tau = E\{ Y(1)  - Y(0) \}.
$$
定义 natural direct effect 为
$$
\NDE = E\{  Y(1,M_0) - Y(0, M_0)  \}. 
$$
定义 natural indirect effect 为
$$
\NIE = E\{  Y(1, M_1) - Y(1, M_0) \} .
$$
\end{definition}


Total effect 是 $Z$ 对 $Y$ 的标准 average causal effect。Natural direct effect 度量的是：若 mediator 被设为没有 intervention 时的 natural value $M_0$，treatment 对 outcome 的效应。Natural indirect effect 度量的是：若 treatment 本身被设为 $z=1$，treatment 通过改变 mediator 所产生的效应。在 composition assumption 下，natural direct and indirect effects 化为
$$
\NDE = E\{  Y(1,M_0) - Y(0)  \} ,\quad
\NIE = E\{  Y(1) - Y(1, M_0) \},
$$
因此，我们可以把 total effect 分解为 natural direct and indirect effects 之和。


\begin{proposition}
\label{prop::mediation-decomposition}
由 \cref{def::mediation} 和 \cref{assume::composition}，有
$$
\tau = \NDE + \NIE .
$$
\end{proposition}

从数学上说，也可以在 treatment 固定为 $0$ 时，把 natural indirect effect 定义为 $ E\{Y(0,M_1)-Y(0,M_0)\}$。不过，这一定义不会导出 \cref{prop::mediation-decomposition} 中的分解。

遗憾的是，由于 interventions 的 cross-world 性质，nested potential outcome $Y(1,M_0)$ 并不是一个容易理解的量：treatment 被设为 $z=1$，但 mediator 却被设为 treatment $z=0$ 下的 natural value $M_0$。显然，对 treatment 的这两个 interventions 不可能在任何一个 realized experiment 中同时发生。为了理解 cross-world potential outcome $Y(1,M_0)$，我们需要想象平行世界的存在，如 \cref{fg::crossworld} 所示。先聚焦于 $Y(1,M_0)$。当 treatment 被设为 $z=1$ 时，mediator 必须取值 $M_1$。如果同时又想把 mediator 设为 $m=M_0$，我们就必须从平行世界中的另一个 experiment 得知同一个 unit 的 $M_0$。这在物理实验上可能并不现实，因为它要求同一个 unit 在两个不同 treatment levels 下同时被干预。在关于 units homogeneity 的某些强 assumptions 下，我们也许可以用另一个 unit 在 control 下的 mediator value 作为 $M_0$ 的 proxy。



 
\subsection{Metaphysics or Science}
\label{sec::popper-science}

Causal inference 本来就困难，甚至连数学记号也没有完全统一。\citet{Robins::1992} 和 \citet{Pearl::2001} 使用 nested potential outcomes 来定义 natural direct and indirect effects。不过，\citet{frangakis2002principal} 把 $Y(1,M_0)$ 和 $Y(0,M_1)$ 称为 {\it a priori counterfactuals}，因为我们无法在任何 physical experiments 中观测到它们。在这个意义上，它们并不先验存在。按照 \citet{popper1963conjectures} 的观点，区分 science 与 metaphysics 的一种方式是 statements 的 falsifiability。也就是说，如果一个 statement 无法基于任何 physical experiments 或 observations 被证伪，那么它就不是 scientific statement，而是 metaphysical statement。由于我们无法在任何 experiments 中观测 $Y(1,M_0)$ 和 $Y(0,M_1)$，除了某些平凡陈述（例如 outcomes 是 binary、continuous 或 bounded）之外，我们无法证伪任何涉及它们的 statements。因此，一个严格的 Popperian statistician 会把 mediation analysis 视为 metaphysics。


更进一步，\citet{dawid2000causal} 批评整个 potential outcomes framework 具有 metaphysical 性质，并把 \cref{sec::po-sciencetable} 中定义的 Rubin's ``Science Table'' 称为 ``metaphysical array''。这不仅是在批评 a priori counterfactuals $Y(1,M_0)$ 与 $Y(0,M_1)$，也在批评更简单的 potential outcomes $Y(1)$ 与 $Y(0)$。\citet{dawid2000causal} 认为，因为我们永远无法联合观测 $Y(1)$ 与 $Y(0)$，所以引入记号 $\{Y(1),Y(0)\}$ 是一种 metaphysical activity。关于 joint distribution $\pr\{Y(1),Y(0)\}$ 的 metaphysical 性质，他是对的；但关于 marginal distributions，他并不正确。基于 observed data，我们确实可以证伪某些关于 marginal distributions 的 statements，尽管我们不能证伪任何关于 joint distribution 的 statements。\footnote{
由 probability theory，给定 marginal distributions $\pr(Y(1)\leq y_1)$ 和 $\pr(Y(0)\leq y_0)$，我们可以用 Frechet--Hoeffding inequality 约束 joint distribution $\pr(Y(1)\leq y_1,Y(0)\leq y_0)$：
\begin{eqnarray*}
&&\max\{  0,  \pr( Y(1) \leq y_1) +\pr( Y(0) \leq y_0 ) - 1  \}  \\
&\leq& 
\pr( Y(1) \leq y_1 , Y(0) \leq y_0 ) \\
&\leq& 
\min\{ \pr( Y(1) \leq y_1) , \pr( Y(0) \leq y_0 )  \}.
\end{eqnarray*}
 这个 inequality 往往较松。遗憾的是，若不施加 additional assumptions，我们没有超出这个 inequality 的任何信息。
} 因此，即使按照 \citet{popper1963conjectures} 的标准，Rubin's Science Table 也不是 metaphysical 的，因为它有一些非平凡且可证伪的 implications，尽管并非所有 implications 都可证伪。这就是 $\{Y(1),Y(0)\}$ 与 $\{Y(1,M_0),Y(0,M_1)\}$ 的根本区别。
 


\begin{figure}[ht]
\centering
\includegraphics[width =  \textwidth]{figures/crossworld.pdf}
\caption{Cross-world potential outcomes $Y(1,M_0)$ and $Y(0, M_1)$}\label{fg::crossworld}
\end{figure}




\section{The Mediation Formula}

\citet{Pearl::2001} 的 mediation formula 依赖下面四个 assumptions。前三个 assumptions 本质上要求：给定 observed covariates 后，treatment 与 mediator 都像是 randomized 的。

\begin{assumption}
\label{assume::mediation1}
不存在 treatment-outcome confounding：
$$
Z\ind Y(z,m)\mid X
$$ 
对所有 $z$ 和 $m$ 成立。
\end{assumption}

\begin{assumption}
\label{assume::mediation2}
不存在 mediator-outcome confounding：
$$
M \ind Y(z,m) \mid (X, Z)
$$ 
对所有 $z$ 和 $m$ 成立。
\end{assumption}

\cref{assume::mediation1} 和 \cref{assume::mediation2} 合在一起通常称为 {\it sequential ignorability}。它们等价于如下 assumption：给定 $X$ 后，$(Z,M)$ jointly randomized：
\begin{eqnarray}\label{eq::joint-randomization}
(Z, M) \ind Y(z,m)\mid X
\end{eqnarray} 
对所有 $z$ 和 $m$ 成立。\cref{eq::joint-randomization} 的证明留作 \cref{para::sr-jr}。



\begin{assumption}
\label{assume::mediation3}
不存在 treatment-mediator confounding：
$$
Z \ind M(z) \mid X
$$ 
对所有 $z$ 成立。
\end{assumption}

最后一个 assumption 是 cross-world independence。

\begin{assumption}
\label{assume::mediation4}
Potential outcomes 与 potential mediators 之间满足 cross-world independence：
$$
Y(z,m) \ind M(z') \mid X
$$ 
对所有 $z,z'$ 和 $m$ 成立。
\end{assumption}


\cref{assume::mediation1,assume::mediation2,assume::mediation3} 都很强，但至少在 treatment 与 mediator randomized 的 experiments 中它们成立。\cref{assume::mediation4} 更强，因为没有 physical experiment 能保证它。若 $z\neq z'$，我们永远无法在任何 experiment 中同时观测 $Y(z,m)$ 与 $M(z')$；因此 \cref{assume::mediation4} 永远无法被验证，从根本上说具有 metaphysical 性质。

下面给出一个例子，其中 \cref{assume::mediation1,assume::mediation2,assume::mediation3,assume::mediation4} 全部成立。

\begin{example}
\label{eg::mediation-assumptions}
给定 $X$，我们生成
\begin{eqnarray*}
Z &=& 1\{  g_Z(X, \varepsilon_Z)  \geq 0 \} ,\\
M(z) &=& 1\{   g_M(X, z, \varepsilon_M)  \geq 0 \} ,\\
Y(z,m) &=& g_Y(X, z, m, \varepsilon_Y ),
\end{eqnarray*}
对 $z,m=0,1$，
其中
$
\varepsilon_Z, \varepsilon_M,
\varepsilon_Y 
$
都是相互独立的 random errors。于是，我们由下式生成 $M$ 和 $Y$ 的 observed values：
\begin{eqnarray*}
M &=& M(Z) = 1\{   g_M(X, Z, \varepsilon_M)  \geq 0 \},\\
Y &=& Y(Z,M) = g_Y(X, Z, M, \varepsilon_Y ) . 
\end{eqnarray*}
可以验证，在这个 data-generating process 下，\cref{assume::mediation1,assume::mediation2,assume::mediation3,assume::mediation4} 成立。
相反，如果允许 $\varepsilon_M$ 和 $\varepsilon_Y$ 分别变为 $\varepsilon_M(z)$ 和 $\varepsilon_Y(z,m)$，那么 \cref{assume::mediation1,assume::mediation2,assume::mediation3,assume::mediation4} 可能失败。
更多细节见 \cref{para::mediation-assumptions-DGP}。
\end{example}



\citet{Pearl::2001} 证明了 mediation analysis 的下面这个关键结果。

\begin{theorem}
\label{thm::mediation}
在 \cref{assume::mediation1,assume::mediation2,assume::mediation3,assume::mediation4} 下，有
$$
E\{ Y(z, M_{z'})  \mid X=x \} =  \sum_m E(Y\mid Z=z, M=m, X=x) \pr(M =  m \mid Z=z', X=x)   
$$
因此，
\begin{eqnarray*}
&&E\{ Y(z, M_{z'})   \}  \\
&=&  \sum_x E\{ Y(z, M_{z'})  \mid X=x \} \pr(X =  x )  \\
&=&  \sum_x \sum_m E(Y\mid Z=z, M=m, X=x) \pr(M =  m \mid Z=z', X=x)   \pr(X =  x ).
\end{eqnarray*}
\end{theorem}

 

\cref{thm::mediation} 假设 $M$ 与 $X$ 都是 discrete 的。对于一般的 $M$ 和 $X$，mediation formulas 变为
$$
E\{ Y(z, M_{z'}) \mid X=x  \} =   \int  E(Y\mid Z=z, M=m, X=x) f ( m \mid Z=z', X=x) \d m
$$
以及
\begin{eqnarray*}
&&E\{ Y(z, M_{z'})   \}  \\
 &=& \int E\{ Y(z, M_{z'}) \mid X=x  \}   f  (x) \d x  \\
 &=& \int \int  E(Y\mid Z=z, M=m, X=x) f ( m \mid Z=z', X=x)     f  (x) \d m \d x   . 
\end{eqnarray*}
由 \cref{thm::mediation} 可知，nested potential outcomes 的均值的 identification formulas 依赖两类量：给定 treatment、mediator 与 covariates 后 outcome 的 conditional mean，以及给定 treatment 与 covariates 后 mediator 的 conditional mean。若 nested potential outcome 涉及 cross-world interventions，我们需要在不同 treatment levels 下评估这两个 conditional means。

如果去掉 cross-world independence assumption，我们可以修改 natural direct and indirect effects 的定义，并且相同的 formulas 仍然成立。更多细节见 \cref{problem::no-crossworld-modification}。

下面给出证明。
 

\begin{myproof}{}{\cref{thm::mediation}}
由 tower property，$E\{Y(z,M_{z'})\}=E[E\{Y(z,M_{z'})\mid X\}]$，所以只需证明 $E\{Y(z,M_{z'})\mid X=x\}$ 的公式。
从 law of total probability 出发，有
\begin{eqnarray*}
&& E\{ Y(z, M_{z'})  \mid X=x\}  \\
&=& 
\sum_m E\{ Y(z, M_{z'})  \mid  M_{z'} = m, X=x\} \pr(  M_{z'} = m\mid X=x) \\
&=&
\sum_m E\{ Y(z, m)  \mid  M_{z'} = m, X=x\} \pr(  M_{z'} = m\mid X=x) \\ 
&=& 
\sum_m \underbrace{ E\{ Y(z, m)  \mid    X=x\} }_{\text{\cref{assume::mediation4}} } 
\underbrace{ \pr(  M  = m\mid Z=z' , X=x) }_{\text{\cref{assume::mediation3}} } \\ 
&=&
\sum_m \underbrace{E( Y  \mid  Z=z, M=m,   X=x)}_{ \text{\cref{assume::mediation1,assume::mediation2}} } 
\pr(  M  = m\mid Z=z' , X=x) .
\end{eqnarray*}
%Averaging over $X$, we obtain the mediation formula. 
\end{myproof}


从数学角度看，上面的证明很直接。它展示了 \cref{assume::mediation1,assume::mediation2,assume::mediation3,assume::mediation4} 的必要性。


条件在 $X=x$ 下，$Y(1,M_1)$ 与 $Y(0,M_0)$ 的 mediation formulas 简化为
\begin{eqnarray*}
&& E\{  Y(1, M_1)\mid X=x \} \\ 
&=& \sum_m E( Y  \mid  Z=1, M=m,   X=x) \pr(M=m \mid Z=1, X=x)\\
&=& E( Y  \mid  Z=1 ,  X=x)
\end{eqnarray*}
以及
\begin{eqnarray*}
&&E\{  Y(0, M_0)\mid X=x \}  \\
&=& \sum_m E( Y  \mid  Z=0, M=m,   X=x) \pr(M=m \mid Z=0, X=x)\\
&=& E( Y  \mid  Z=0 ,  X=x)
\end{eqnarray*}
这是基于 law of total probability；而 $Y(1,M_0)$ 的 mediation formula 简化为
\begin{eqnarray*}
&&E\{  Y(1, M_0)\mid X=x \}  \\
&=&  \sum_m E( Y  \mid  Z=1, M=m,   X=x) \pr(M=m \mid Z=0, X=x),
\end{eqnarray*}
其中 outcome 的 conditional expectation 是在 $Z=1$ 下给出的，但 mediator 的 conditional distribution 是在 $Z=0$ 下给出的。
这导出了 natural direct and indirect effects 的 identification formulas。

\begin{corollary}
\label{corollary::mediation-nde-nie}
在 \cref{assume::mediation1,assume::mediation2,assume::mediation3,assume::mediation4} 下，
conditional natural direct and indirect effects 可由下式识别：
\begin{eqnarray*}
\NDE(x) &=& E\{ Y(1, M_0) - Y(0, M_0) \mid X=x \} \\
&=& \sum_m \{ E( Y  \mid  Z=1, M=m,   X=x) - E( Y  \mid  Z=0, M=m,   X=x) \} \\
&& \qquad  \times \pr(M=m \mid Z=0, X=x)
\end{eqnarray*}
以及
\begin{eqnarray*}
\NIE(x) &=& E\{ Y(1, M_1) - Y(1, M_0) \mid X=x \} \\
&=& \sum_m   E( Y  \mid  Z=1, M=m,   X=x)\\
&&  \qquad  \times  \{ \pr(M=m \mid Z=1, X=x) -  \pr(M=m \mid Z=0, X=x) \};
\end{eqnarray*}
unconditional effects 可由 $\NDE=\sum_x \NDE(x)\pr(X=x)$ 和 $\NIE=\sum_x \NIE(x)\pr(X=x)$ 识别。
\end{corollary}





作为一个 special case，当 $M$ 为 binary 时，$\NIE$ 的公式化为下面的 product form。


\begin{corollary}
\label{corollary::mediation-binary-M}
在 \cref{assume::mediation1,assume::mediation2,assume::mediation3,assume::mediation4} 下，若 mediator $M$ 为 binary，则有
$$
\NIE(x)  =      \tau_{Z\rightarrow M} (x) \tau_{M\rightarrow Y} (1,x)
$$
以及
$
\NIE = E\{  \NIE(X)   \},
$
其中
$$
\tau_{Z\rightarrow M}(x) =  \pr(M=1 \mid Z=1, X=x) -  \pr(M=1 \mid Z=0, X=x) . 
$$
以及
$$
\tau_{M\rightarrow Y} (z,x) = E( Y  \mid  Z=z, M=1,   X=x) - E( Y  \mid  Z=z, M=0,   X=x)
$$
\end{corollary}

\cref{corollary::mediation-binary-M} 的证明留作 \cref{para::nie-binary-m}。
\cref{corollary::mediation-binary-M} 给出了 binary $M$ 情形下的一个简单公式。
当给定 $X$ 后 $Z$ randomized 时，可以把 $\tau_{Z\rightarrow M}(x)$ 看作 $Z$ 对 $M$ 的 conditional average causal effect。当给定 $(X,Z)$ 后 $M$ randomized 时，可以把 $\tau_{M\rightarrow Y}(z,x)$ 看作 $M$ 对 $Y$ 的 conditional average causal effect。Conditional natural indirect effect 等于二者的乘积。这与我们的直觉一致：indirect effect 先从 $Z$ 作用到 $M$，再从 $M$ 作用到 $Y$。


\section{The Mediation Formula Under Linear Models}


\cref{thm::mediation} 给出了 mediation analysis 的 nonparametric identification formula。它允许我们在不同 models 下推导各种 mediation analysis formulas。下面我会介绍 linear models 下著名的 Baron--Kenny method。
\citet{vanderweele2015explanation} 为许多常用 models 给出了 natural direct and indirect effects 的 explicit formulas。其它 models 的细节留到 \cref{chapter::mediation-problmes}。

\subsection{The Baron--Kenny Method}


\begin{figure}[ht]
\centering
$$
\begin{xy}
\xymatrix{
&  & X \ar[ddll] \ar[dd]^{\beta_2} \ar[ddrr]^{\theta_4} \\
&&\\
Z \ar[rr]^{\beta_1} \ar@/_2.8pc/[rrrr]_{\theta_1} & & M\ar[rr]^{\theta_2} & & Y & \text{indirect effect: } \beta_1\theta_2 \\
&&&&& \text{  direct effect: } \theta_1 }
\end{xy}
$$
\caption{Linear models 下 mediation 的 Baron--Kenny method}\label{fg::mediationDAG-BK}
\end{figure}

Baron--Kenny method 假设：给定 treatment 与 covariates 后，mediator 与 outcome 满足下面的 linear models。

\begin{assumption}
[linear models for the Baron--Kenny method]
\label{assume::mediation-bk-model}
Mediator 与 outcome 都服从 linear models：
$$
\left\{\begin{array}{rcl}
E(M \mid Z, X) &=&  \beta_0 + \beta_1 Z + \beta_2\tran X, \\
E(Y\mid Z,M,X) &=& \theta_0 + \theta_1 Z + \theta_2 M + \theta_4\tran X.
\end{array}\right.
$$
\end{assumption}


 
在这些 linear models 下，natural direct and indirect effects 的公式简化为 coefficients 的函数。

\begin{corollary}
[Baron--Kenny formulas for mediation]
\label{corollary::bk-formulas}
在 \cref{assume::mediation1,assume::mediation2,assume::mediation3,assume::mediation4,assume::mediation-bk-model} 下，有
$$
\NDE = \theta_1
,\quad 
\NIE = \theta_2 \beta_1 .
$$
\end{corollary}


\cref{corollary::bk-formulas} 中的公式可以从 \cref{fg::mediationDAG-BK} 直观理解。Direct effect 等于 path $Z\rightarrow Y$ 上的 coefficient，而 indirect effect 等于 path $Z\rightarrow M\rightarrow Y$ 上 coefficients 的乘积。下面基于 \cref{thm::mediation} 给出证明。


\begin{myproof}{}{\cref{corollary::bk-formulas}}
Conditional NDE 等于
$$
\NDE(x)  
= \sum_m  \theta_1  \pr(M=m \mid Z=0, X=x) = \theta_1
$$
而 conditional NIE 等于
\begin{eqnarray*}
\NIE(x) &=&  \sum_m  (\theta_0 + \theta_1  + \theta_2 m + \theta_4\tran x )\\
&&  \qquad  \times  \{ \pr(M=m \mid Z=1, X=x) -  \pr(M=m \mid Z=0, X=x) \} \\
&=& \theta_2 \{ E(M=m \mid Z=1, X=x) -  E(M=m \mid Z=0, X=x) \} \\
&=&  \theta_2 \beta_1 ,
\end{eqnarray*}
它们不依赖于 $x$。因此，它们也是 unconditional natural direct and indirect effects 的公式。
\end{myproof}


 
如果得到这些 coefficients 的 OLS estimators，则可以用
$$
\hat{\NDE} = \hat{\theta}_1,\qquad
\hat{\NIE} = \hat{\theta}_2 \hat{\beta}_1,
$$
来估计 direct and indirect effects。这称为 Baron--Kenny method \citep{Judd::1981, baron1986moderator}，尽管它已有若干 antecedents \citep[e.g.,][]{hyman1955survey, alwin1975decomposition, Judd::1981, sobel1982asymptotic}。


Standard software packages 会从 OLS 中报告 $\hat{\NDE}$ 的 standard error。\citet{sobel1982asymptotic, sobel1986some} 使用 delta method 得到 $\hat{\NIE}$ 的 standard error。
%Based on the heuristics:
%\begin{eqnarray*}
%\hat{\theta}_2  \hat{\beta}_1 -  \theta_2 \beta_1
%&=& \hat{\theta}_2  \hat{\beta}_1 -  \theta_2 \hat{\beta}_1 +   \theta_2 \hat{\beta}_1  -  \theta_2 \beta_1 \\
%&=&  (\hat{\theta}_2 -  \theta_2 )  \hat{\beta}_1 + \theta_2 (\hat{\beta}_1 - \beta_1 ) \\
%&\approx& (\hat{\theta}_2 -  \theta_2 )   \beta_1 + \theta_2 (\hat{\beta}_1 - \beta_1 )
%\end{eqnarray*}
%and the independence of $\hat{\theta}_2 $ and $ \hat{\beta}_1$, 
由 \cref{example::asymptotic-normal-product} 中的公式，$\hat{\theta}_2\hat{\beta}_1$ 的 asymptotic variance 等于
$
\var(\hat{\theta}_2)  \beta_1^2 + \theta_2^2 \var(\hat{\beta}_1) .
$
因此 estimated variance 为
$$
\hat{\var}(\hat{\theta}_2)  \hat{\beta}_1^2 + \hat{\theta}_2^2 \hat{\var}(\hat{\beta}_1) .
$$
在 mediation analysis 文献中，基于 $\hat{\theta}_2\hat{\beta}_1$ 以及上面的 estimated variance 检验 $\NIE$ 的 null hypothesis，称为 {\it Sobel's test}。


 


\subsection{An Example}\label{section::BK-jobs}

可以很容易地用下面的代码实现 Baron--Kenny method。

 \begin{lstlisting}
library("car")
BKmediation = function(Z, M, Y, X)
{
  ## two regressions and coefficients
  mediator.reg   = lm(M ~ Z + X)
  mediator.Zcoef = mediator.reg$coef[2]
  mediator.Zse   = sqrt(hccm(mediator.reg)[2, 2])
  
  outcome.reg    = lm(Y ~ Z + M + X)
  outcome.Zcoef  = outcome.reg$coef[2]
  outcome.Zse    = sqrt(hccm(outcome.reg)[2, 2])
  outcome.Mcoef  = outcome.reg$coef[3]
  outcome.Mse    = sqrt(hccm(outcome.reg)[3, 3])
  
  ## Baron-Kenny point estimates
  NDE = outcome.Zcoef
  NIE = outcome.Mcoef*mediator.Zcoef
  
  ## Sobel's variance estimate based the delta method
  NDE.se = outcome.Zse
  NIE.se = sqrt(outcome.Mse^2*mediator.Zcoef^2 + 
                  outcome.Mcoef^2*mediator.Zse^2)
  
  res = matrix(c(NDE, NIE, 
                 NDE.se, NIE.se,
                 NDE/NDE.se, NIE/NIE.se), 
               2, 3)
  rownames(res) = c("NDE", "NIE")
  colnames(res) = c("est", "se", "t")
  
  res
}
 \end{lstlisting}
 
 
 回到 \cref{eg::mediation-jobs}，可以得到 direct and indirect effects 的如下 estimates：
 
 \begin{lstlisting}
> jobsdata = read.csv("jobsdata.csv")
> Z = jobsdata$treat
> M = jobsdata$job_seek
> Y = jobsdata$depress2
> getX     = lm(treat ~ econ_hard + depress1 + 
+                 sex + age + occp + marital + 
+                 nonwhite + educ + income,
+               data = jobsdata)
> X = model.matrix(getX)[, -1]
> res = BKmediation(Z, M, Y, X)
> round(res, 3)
       est    se      t
NDE -0.037 0.042 -0.885
NIE -0.014 0.009 -1.528
\end{lstlisting}

Direct and indirect effects 的 estimates 都为负，尽管它们并不显著。



\section{Sensitivity analysis}


Mediation analysis 依赖强且不可检验的 assumptions。一个关键 assumption 是 treatment、mediator 与 outcome 之间不存在 unmeasured confounding。文献中已有多种 sensitivity analysis methods。特别地，\citet{ding2016sharp} 提出了 Cornfield-type sensitivity bounds，\citet{zhang2022interpretable} 基于 linear structural equation models 提出了一种专门针对 Baron--Kenny method 的 sensitivity analysis method。这些内容超出了本书范围。





 

\section{Homework problems}\label{chapter::mediation-problmes}



\homeworksolutionref{27}
\paragraph{\citet{Robins::1992}'s clarification of the units based on nested potential outcomes}\label{hw::robins1992-latentgroups}


假设 $(Z,M,Y)$ 都是 binary。根据
$$
M(1), M(0),  Y(1, M_1), Y(1, M_0), Y(0, M_1), Y(0, M_0) ,
$$
的 joint value，我们有多少种 units 类型？

如果假设 $Z$ 对 $M$ 的 monotonicity，即
$$
M(1) \geq  M(0),
$$
那么有多少种 units 类型？

如果进一步假设 $Z$ 和 $M$ 对 $Y$ 的 monotonicity，即
$$
Y(1, m) \geq Y(0, m),  \quad Y(z, 1) \geq Y(z, 0)
$$
对所有 $z$ 和 $m$ 成立，那么有多少种 units 类型？


\paragraph{Sequential randomization and joint randomization}\label{para::sr-jr}

证明 \cref{eq::joint-randomization} 等价于 \cref{assume::mediation1,assume::mediation2}。


\paragraph{Verifying the assumptions for mediation analysis}\label{para::mediation-assumptions-DGP}

证明在 \cref{eg::mediation-assumptions} 的 data generating process 下，\cref{assume::mediation1,assume::mediation2,assume::mediation3,assume::mediation4} 成立。
还要证明，如果允许 $\varepsilon_M$ 和 $\varepsilon_Y$ 分别变为 $\varepsilon_M(z)$ 和 $\varepsilon_Y(z,m)$，那么 \cref{assume::mediation1,assume::mediation2,assume::mediation3,assume::mediation4} 可能失败。
 
 \paragraph{Another set of assumptions for the mediation formula}\label{para::alternative-si}



\citet{Imai::2010} 使用下面这组 assumptions 推导 mediation formula。

\begin{assumption}
\label{assume::mediation-imai}
有
$$
\{Y(z, m), M(z') \} \ind  Z\mid X
$$
以及
$$
Y(z, m) \ind M(z' ) \mid (Z=z' , X)
$$
对所有 $z,z',m$ 成立。
\end{assumption}


\begin{theorem}
\label{thm::mediation-imai}
在 \cref{assume::mediation-imai} 下，\cref{thm::mediation} 中的 mediation formula 成立。
\end{theorem}


证明 \cref{thm::mediation-imai}。



\paragraph{Difference method and product method are identical}\label{hw::difference=product}


在正文中，我们首先得到 OLS fits
$$
\left\{\begin{array}{rcl}
\hat M_i &=&  \hat\beta_0 + \hat\beta_1 Z_i + \hat\beta_2\tran X_i, \\
\hat Y_i &=&\hat \theta_0 + \hat\theta_1 Z_i +\hat \theta_2 M_i + \hat\theta_4\tran X_i.
\end{array}\right.
$$
NIE 的 estimate 是乘积 $\hat\theta_2\hat\beta_1$，这称为 product method。

或者，也可以先得到 OLS fits：
$$
\left\{\begin{array}{rcl}
\hat Y_i &=&  \hat\alpha_0 + \hat\alpha_1 Z_i + \hat\alpha_1 \tran X_i, \\
\hat Y_i &=&\hat \theta_0 + \hat\theta_1 Z_i +\hat \theta_2 M_i + \hat\theta_4\tran X_i.
\end{array}\right.
$$
NIE 的另一个 estimate 是差 $\hat\alpha_1-\hat\theta_1$，这称为 difference method。

证明 $\hat\alpha_1-\hat\theta_1=\hat\theta_2\hat\beta_1$。


Remark: 回顾 \cref{hw::cochran+ovb} 中的 Cochran's formula。



\paragraph{Natural indirect effect with a binary mediator}\label{para::nie-binary-m}

证明 \cref{corollary::mediation-binary-M}。


\paragraph{With treatment-outcome interaction on the outcome}\label{para::interaction}
\citet{vanderweele2015explanation} 建议使用下面的 linear models：
$$
\left\{\begin{array}{rcl}
E(M \mid Z, X) &=&  \beta_0 + \beta_1 Z + \beta_2\tran X, \\
E(Y\mid Z,M,X) &=& \theta_0 + \theta_1 Z + \theta_2 M + \theta_3 ZM + \theta_4\tran X,
\end{array}\right.
$$
其中 outcome model 包含 treatment 与 mediator 之间的 interaction term。


在上述 linear models 下，证明
$$
\NDE = \theta_1 + \theta_3\{ \beta_0 + \beta_2\tran E(X)\}, \qquad 
\NIE = (\theta_2 + \theta_3 ) \beta_1. 
$$
如何用 IID data 估计 $\NDE$ 与 $\NIE$？

Remark:
考虑 $Z$ 与 $M$ 都为 binary 的简单情形。
在线性模型下，$Z$ 对 $M$ 的 average causal effect 等于 $\beta_1$，$M$ 对 $Y$ 的 average causal effect 等于 $\theta_2+\theta_3E(Z)$。因此，可能这两个 effects 都为正，但 natural indirect effect 为负。例如：
$$
\beta_1 = 1,\quad \theta_2 = 1, \quad  \theta_3= -1.5, \quad E(Z) = 0.5.
$$ 
这有些 paradoxical，可称为 {\it mediator paradox}。\citet{chen2007criteria} 报告了一个相关的 {\it surrogate endpoint paradox} 或 {\it intermediate variable paradox}。



\paragraph{Mediation analysis with continuous mediator and binary outcome}\label{para::continM-binaryY}


考虑下面 mediator 的 Normal linear model 与 binary outcome 的 logistic model：
$$
\left\{\begin{array}{rcl}
M \mid Z, X &\sim &  \textsc{N}(\beta_0 + \beta_1 Z + \beta_2\tran X, \sigma_M^2), \\
\logit \{ \pr(Y = 1 \mid Z,M,X ) \} &=& \theta_0 + \theta_1 Z + \theta_2 M +   \theta_4\tran X,
\end{array}\right.
$$
其中 $\logit(w)=\log\{w/(1-w)\}$，其 inverse 为 $\expit(w)=(1+e^{-w})^{-1}$。
用 model parameters 和 $X$ 的分布表示 $\NDE$ 与 $\NIE$。如何用 IID data 估计 $\NDE$ 与 $\NIE$？




\paragraph{Mediation analysis with binary mediator and continuous outcome}\label{para::logit-binary-m}
考虑下面 binary mediator 的 logistic model 与 outcome 的 linear model：
$$
\left\{\begin{array}{rcl}
\logit\{ \pr(M  = 1\mid Z, X) \}&=&  \beta_0 + \beta_1 Z + \beta_2\tran X, \\
E(Y\mid Z,M,X) &=& \theta_0 + \theta_1 Z + \theta_2 M +   \theta_4\tran X . 
\end{array}\right.
$$
%where $\logit (w)  = \log\{ w/(1-w) \} $ with inverse $\expit(w) = (1 + e^{-w})^{-1}.$


在这些 models 下，证明
$$
\NDE = \theta_1  ,\qquad
\NIE = \theta_2 E \left\{  \expit(\beta_0 + \beta_1  + \beta_2\tran X)
-  \expit(\beta_0   + \beta_2\tran X)  \right\}. 
$$
如何用 IID data 估计 $\NDE$ 与 $\NIE$？






\paragraph{Mediation analysis with binary mediator and outcome}
\label{para::mediation-binaryMY}

考虑下面 binary mediator 与 binary outcome 的 logistic models：
$$
\left\{\begin{array}{rcl}
\logit\{ \pr(M  = 1\mid Z, X) \}&=&  \beta_0 + \beta_1 Z + \beta_2\tran X, \\
\logit \{ \pr(Y = 1 \mid Z,M,X ) \} &=& \theta_0 + \theta_1 Z + \theta_2 M +   \theta_4\tran X.
\end{array}\right.
$$
用 model parameters 和 $X$ 的分布表示 $\NDE$ 与 $\NIE$。如何用 IID data 估计 $\NDE$ 与 $\NIE$？



  
\paragraph{Modify the definitions to drop the cross-world independence}\label{problem::no-crossworld-modification}


定义
$$
Y(z, F_{M_{z'}\mid X}) = \int Y(z,m) f(M_{z'} = m\mid X) \d m
$$
为 treatment $z$ 下、并从 $M_{z'}\mid X$ 的分布中随机抽取 mediator 时的 potential outcome。
当 $M$ 为 discrete 时，定义简化为
$$
Y(z, F_{M_{z'}\mid X}) =\sum_m Y(z,m) \pr(M_{z'} = m\mid X).
$$
$Y(z,M_{z'})$ 与 $Y(z,F_{M_{z'}\mid X})$ 的关键区别在于：$M_{z'}$ 是同一个 unit 的 potential mediator，而 $F_{M_{z'}\mid X}$ 是从 whole population 中 potential mediator 的 conditional distribution 进行的一次 random draw。
定义 natural direct and indirect effects 为
$$
\NDE = E\{Y(1, F_{M_{0}\mid X})  - Y(0, F_{M_{0}\mid X})   \},\quad
\NIE  =  E\{Y(1, F_{M_{1}\mid X})  - Y(1, F_{M_{0}\mid X})   \}.
$$


证明在 \cref{assume::mediation1,assume::mediation2,assume::mediation3} 下，$\NDE$ 与 $\NIE$ 的 identification formulas 仍与正文相同。



Remark: 修改 nested potential outcomes 的定义可以让我们放松很强的 cross-world independence assumption，但会削弱 natural direct and indirect effects 的解释。更多讨论见 \citet{vanderweele2015explanation}；关于 time-varying treatment and mediator 这一更复杂 setting 中的应用，见 \citet{vanderweele2017mediation}。

  
  
 
